\section{Team}

\noindent \textbf{Project Tutor:} dr inż. Aleksandra Karpus

\noindent \textbf{Project Authors:}
\begin{itemize}
   \item Wiktor Czetyrbok
   \item Kamil Dędza
   \item Agnieszka Kuleta
\end{itemize}

\subsection{Organization of the Team}

\begin{table}[h!]
\centering
\begin{tabular}{| m{2cm} | m{2cm} | m{4cm} | m{3cm} | m{2cm} |}
    \hline
    \textbf{Team Member} & \textbf{Role} & \textbf{Experience} & \textbf{Skills} & \textbf{Collaboration Type} \\
    \hline
    Kamil Dędza & Software Developer & Junior Data Engineer at DXC & Python, PL/SQL, JAVA & Remotely \\
    \hline
    Wiktor Czetyrbok & Software Developer & Software Engineer at GridDynamics & JAVA, JavaScript, SQL, GCP, DevOps & Remotely \\
    \hline
    Agnieszka Kuleta & Scrum Master & Data Analyst at Nordea Bank & MS SQL, Python, JAVA & Remotely \\
    \hline
\end{tabular}
\caption{Team}
\end{table}

\subsection{Project Responsibilities}
\begin{itemize}
   \item \textbf{Wiktor Czetyrbok:} Mainly responsible for application architecture and design.
   \item \textbf{Kamil Dędza:} Mainly responsible for data management and integration.
   \item \textbf{Agnieszka Kuleta:} Mainly responsible for business requirements, communication with the product owner, and application design.
\end{itemize}

\subsection{Communication Within the Team}
\begin{itemize}
    \item \textbf{Communication Methods:} We communicate via email, chat, and video calls.
    \item \textbf{Communication Tools:} Messenger, Discord, Gdańsk University of Technology's Mail.
    \item \textbf{Meetings:} We meet at least twice a week, primarily using Discord.
    \item \textbf{Communication with Project Environment:} Every two weeks on Discord on Fridays at 2 PM (meetings with the product owner/tutor).
\end{itemize}

\subsection{Sharing Code and Documents}
\begin{itemize}
    \item \textbf{Code:} Available on \href{https://github.com/culinary-portal}{GitHub} (https://github.com/culinary-portal).
    \item \textbf{Documents:} Available on \href{https://drive.google.com/drive/folders/1RUclEIFVWzzO89Ef-BEOq56EYikbE4ET?usp=drive_link}{Google Drive}.
\end{itemize}

\section{Scrum}

 Scrum is an Agile methodology that though its iterative progress of  short cycles \textbf{sprints} which last 2 weeks. It allows us to build new features, such as recipe database, user accounts, and interactive interfaces, while continually gathering feedback from stakeholders. It also enables us to respond to changing requirements, making sure that final product is closer to future users needs. This method also increases effectiveness of team communication and helps to prioritize tasks based on project goals. 
 
 \subsection{Backlog Documentation}

After creating a high level overview of all the requirements \ref{chapter:requirements}, there is a necessity to introduce technical tasks. They are in fact a bridge between business requirements and detailed work of the developers. They answer the question how the team can achieve particular features, and how to implement them. Some stakeholders at first glance might not notice them, but they are crucial during the process of creating a solution. These types of product backlog items are also referred to as “technical debt” because, much like financial debt, they can cause problems if allowed to build up. It is wise to prioritize technical work alongside features and defects' \cite{TechnicalTaskBacklog}. These tasks are usually written in clear language, what must be done under the hood, for example configuring the database or creating a communication between different services. They also help with complex problems, dividing them into smaller, separable chunks. The team of developers can in a clear way distribute the work between the members. Together, with previously mentioned business requirements they present the whole amount of work and resources needed to achieve the success and to create a solution. Planning can be also done only if the majority of requirements are gathered, and only then it is possible to estimate the time and work allocation. Additionally, this process offers individuals without extensive technical knowledge a clearer understanding of the complexity of certain features and the variety of tasks they involve, as outlined in  at \textbf{Appendix A} \ref{chapter:appendixA}.
\par
The project backlog was created and managed using \textbf{Jira}, a software that enables the definition and prioritization of tasks. Each backlog item, such as epics, user stories, and tasks, was added to the system, providing easy access to detailed descriptions, statuses, and time estimates. In line with the scrum methodology, the organization of sprints played a crucial role in ensuring continuous progress, adaptability to changes, and frequent delivery of valuable increments. Every two weeks, our team defined goals and selected previously prioritized items based on the time and complexity of tasks, maintaining a structured and efficient workflow. 


\section{Tools}
\subsection{Supporting tools}
As supporting tools for areas such as backlog management, system development, and testing, we utilized a combination of platforms. Under the subsection of Work Organization, we incorporated tools such as\textbf{Jira} was used for tasks and backlog management, allowing the team to define, prioritize, and track tasks efficiently. For version control and collaborative development, we used \textbf{GitHub} and to host our services we have used \textbf{Google Cloud Platform}. \textbf{Google Drive} served as a centralized platform for storing and sharing project documentation, ensuring accessibility and organization of resources. Additionally, \textbf{Discord} was our primary communication tool for real-time discussions and coordination, complemented by \textbf{Email} via Outlook for formal communication. As a Programming tools we have used:

\begin{itemize}
    \item Code Editors:
        \begin{itemize}
            \item Visual Studio Code
            \item IntelliJ IDEA
            \item PyCharm
        \end{itemize}
        \item Version Control git:
        \begin{itemize}
            \item Git
            \item Github
        \end{itemize}
        \item Dependecy and Package Managment:
        \begin{itemize}
            \item pip (Pip Installs Packages)
            \item Gradle (Java)
            \item npm (Node Package Manager)
        \end{itemize}
        \item Database and management tools:
        \begin{itemize}
            \item PostgreSQL
        \end{itemize}
        \item {Libraries and Frameworks:}
\begin{itemize}
    \item Angular   
    \item Spring boot (Java)
\end{itemize}
\end{itemize}

\begin{itemize}
    \item Programming Languages:
        \begin{itemize}
            \item Python
            \item Java
            \item JavaScript
            \item HTML, SCSS
        \end{itemize}
    \item Web frameworks:
        \begin{itemize}
            \item Spring (Java)
        \end{itemize}
    \item Database:
        \begin{itemize}
            \item PostgreSQL/Cloud SQL
            \item Redis/Memorystore
        \end{itemize}
    \item Frontend Technologies:
        \begin{itemize}
            \item JavaScript
        \end{itemize}
    \item Libraries and tools for building user interface: 
        \begin{itemize}
            \item Angular and Tailwind
    \end{itemize}
\end{itemize}



